% !TEX root = ../main.tex
% ---------------------------------------------------------------------
%  P5 -- Multicore
% ---------------------------------------------------------------------
\vlsec{cP5}{P5}{Multicore}
\label{sec:multicore}

\begin{formel}[Amdahl für Kerne]
\[S(N)=\frac{1}{(1-p)+\dfrac{p}{N}}\]
$p=80\,\%$ parallel, $N=2$:
$S=\frac{1}{0.2+0.4}=1.67$. Dieselbe Formel wie für SIMD-Lanes
(Abschnitt~\ref{form:amdahl}) -- nur heisst $s$ jetzt $N$.
\end{formel}

\SW{Gründe für mehrere Kerne}
\begin{itemize}
\item Komplexität von HW nach SW verschieben: mehrere kleine schnelle
      RISC-Kerne statt eines grossen superskalaren Out-of-Order-Kerns
      -- aber \emph{zuerst} die Caches vergrössern, falls das Optimum
      noch nicht erreicht ist
\item Lizenzkosten steigen typisch \emph{nicht} linear mit der Kernzahl
\item little-big-Architekturen sparen Strom (z.\,B.\ 2 Effizienz- $+$
      2 Performance-Kerne)
\end{itemize}
\SW{HW-Grenze} RAM-Bandbreite. \SW{Nachteil} Datenkohärenz.

\subsection{Multitasking-OS (FreeRTOS)}
Pro Prozessor läuft genau \emph{ein} Task. Bei jedem System-Tick (z.\,B.\
20\,ms Timer-IRQ) kann der Scheduler nach Priorität wechseln -- technisch
nur durch Austausch des \T{Stackpointers}; jeder Task hat seinen eigenen
Stack. Neue Tasks aus der \emph{ready}-Liste; wartende wandern in die
\emph{blocked}-Liste.
\SW{Vorteil} keine Cycles für Polling, einfache Programmierung.
\SW{Problem} Datenkohärenz und geteilte Ressourcen.

\begin{falle}[Reentranz ist Pflicht]
Nicht reentrant: \R{temp} als \emph{globale} Variable in \R{swap()} --
zwei gleichzeitig aktive Instanzen zerstören sich gegenseitig.
Reentrant: \R{temp} \emph{lokal} (auf dem eigenen Stack) -- dann dürfen
Tasks und Kerne denselben Code teilen.
\end{falle}

\subsection{Sechs Ursachen für Kohärenzprobleme}
\label{sec:kohaerenz}
\begin{ztab}{@{}lX@{}}
Stale Data & Core 1 liest den alten Wert aus seinem Cache \\
Write-Write & beide schreiben gleichzeitig, Ergebnis unbestimmt \\
Read-Modify-Write & beide lesen 10, schreiben 11 -- ein Inkrement weg \\
Non-Atomic Multiword & 64\,Bit als zwei 32-Bit-Writes, Leser sieht Mix \\
Out-of-Order & \R{flag=1} wird vor \R{data=42} sichtbar \\
DMA vs.\ CPU & DMA liest, während die CPU noch schreibt \\
\end{ztab}

\begin{bsp}[Race Condition: \R{counter++}]
\begin{lstlisting}
counter++  ->  LOAD counter / ADD 1 / STORE counter

Core0: LOAD 42      Core1: LOAD 42
       ADD 1               ADD 1
       STORE 43            STORE 43   -> 43 statt 44
\end{lstlisting}
Je $10^6$ Inkremente ergeben \emph{nicht} $2\cdot10^6$.
Lösung: \R{lock(); counter++; unlock();} bzw.\ konzeptionell
\R{atomic\_fetch\_add(\&counter,1)}.
\end{bsp}

\subsection{Kooperation paralleler Prozesse}
\SW{Mutual Exclusion brutal} Scheduler, andere CPUs und Interrupts
abschalten und Caches flushen. \SW{Sauber} Semaphore/Mutex.
\begin{ztab}{@{}lX@{}}
Kommunikation & \T{Mailbox} (blockiert) $|$ \T{Message Queue} \\
Synchronisation & \T{Semaphore} (zählend, z.\,B.\ 3 gleiche Ressourcen)
  $|$ \T{Mutex} $|$ \T{Flags} \\
\end{ztab}

\begin{falle}[Priority Inversion]
Ein niederpriorer Task hält das Semaphor und blockiert damit einen
hochprioren; ein mittelpriorer verdrängt den niederprioren dauernd
\ra der hochpriore wartet indirekt auf einen Task, der nie drankommt.\\
\T{Lösung:} \T{Mutex} mit \emph{Priority Inheritance} -- der Halter erbt
vorübergehend die hohe Priorität.\\
\emph{Mars Pathfinder 1997}: genau dieser Fehler, sporadische Resets;
Fix war Priority Inheritance in VxWorks.
\end{falle}

\SW{SMP} identische CPUs, \emph{ein} OS.
\SW{AMP} eine OS-Instanz pro Kern, braucht HW-Unterstützung für
Semaphore/Mutex.

\subsection{FreeRTOS-Kurzreferenz}
\begin{ztab}{@{}lX@{}}
\R{xTaskCreate(fn,name,stack,par,prio,\&h)} & Task anlegen \\
\R{vTaskStartScheduler()} & Scheduler starten \\
\R{vTaskDelay(ms)} & blockieren (nicht pollen!) \\
\R{xQueueCreate(len,size)} & Queue anlegen \\
\R{xQueueSend/Receive(q,\&v,wait)} & Rückgabe \R{pdPASS} prüfen \\
\R{xMessageBufferSend/Receive} & variable Nachrichtenlänge \\
\R{xSemaphoreCreateMutex()} & Mutex (mit Prio-Vererbung) \\
\R{xSemaphoreTake(m,portMAX\_DELAY)} & kritischen Abschnitt betreten \\
\R{xSemaphoreGive(m)} & verlassen \\
\R{xTaskNotifyWait} / \R{\dots FromISR} & leichtgewichtiges Signal \\
\R{vTaskCoreAffinitySet(h,1<<n)} & Task an Kern $n$ binden \\
\end{ztab}
SMP-Konfiguration: \R{configNUM\_CORES 2},
\R{configRUN\_MULTIPLE\_PRIORITIES}, \R{configUSE\_CORE\_AFFINITY},
\R{configUSE\_PREEMPTION}.

\subsection{RP2350 Multiprozessor-Subsystem}
$2\times$Armv8-M (DSP, FPU, Security) \emph{oder} $2\times$\R{RV32IMAC};
Mischbetrieb über \R{ARCHSEL} möglich, braucht aber zwei Programmimages.
\begin{itemize}
\item symmetrische Bussicht: beide Kerne erreichen alle Ressourcen gleich
      schnell
\item pro Kern zwei 32-Bit-Links zum Systembus (Instruktion $+$
      Load/Store $=$ Harvard)
\item ungenutzte Kerne in Reset mit abgeschaltetem Takt
      \ra \emph{null} dynamische Leistung (CMOS)
\item alle Interrupts gehen an beide Kerne, der interne Controller abonniert
\end{itemize}

\SW{SIO (Single-cycle IO)}
\begin{ztab}{@{}lX@{}}
Spinlocks & 32 Stück, Ein-Bit-Flag je Adresse. \emph{Lesen} $=$ Anspruch
  (Rückgabe $\neq0$ \ra erhalten), \emph{Schreiben} (beliebiger Wert) $=$
  freigeben. Bei Gleichstand gewinnt Core\,0. Nur für sehr kurze
  kritische Abschnitte (Spinning ist teuer). \\
FIFOs & Inter-Processor Mailboxes, 32\,Bit breit, 4 Einträge tief;
  Flags \R{VLD}/\R{RDY}/\R{ROE}/\R{WOF} (letzte zwei sticky); IRQ\,25 \\
Doorbells & 8 Flags je Richtung, ein IRQ pro Kern (26), für
  \emph{kumulative, ungeordnete} Ereignisse \\
Rest & \R{CPUID} (0/1), Integer-Divider, Interpolator,
  RISC-V-Platform-Timer, GPIO fürs Bitbanging; alles doppelt für
  secure/non-secure \\
\end{ztab}

\begin{formel}[Global Exclusive Monitor (empfohlen statt Spinlocks)]
Processing Elements: Core\,0, Core\,1 und \emph{der DMA}.
Reservation Page $=$ beliebige 16-B-ausgerichtete SRAM-Region.
Ein exclusive Write gelingt nur, wenn
(a)~ein exclusive Read desselben PE voranging,
(b)~nur eine Reservation offen ist,
(c)~Adresse, Grösse und Rechte übereinstimmen und
(d)~kein anderes PE seither in dieselbe Granule geschrieben hat.
Sonst wird der Write verworfen und meldet das zurück.
\end{formel}
\begin{lstlisting}
arm_acquire:               lr.w rd,(rs1)
  LDREX   r1,[r0]            ; Reservation setzen
  CMP     r1,#0            sc.w rd,rs2,(rs1)
  STREXEQ r1,r2,[r0]         ; rd=0 ok, !=0 fehlgeschlagen
  CMPEQ   r1,#0
  BNE     acquire
\end{lstlisting}
RISC-V zusätzlich \R{amoadd/amoswap/amoand/\dots} (Extension \R{A}).
